\documentclass[12pt]{article}
\usepackage[margin=1in]{geometry}
\usepackage{times}
\usepackage[hidelinks]{hyperref}
\usepackage{titlesec}
\usepackage{enumitem}

\titleformat{\section}{\normalfont\bfseries\large}{}{0pt}{}
\titlespacing{\section}{0pt}{10pt}{4pt}
\pagenumbering{gobble}
\setlength{\parindent}{0pt}
\setlength{\parskip}{3pt}
\linespread{0.97}

\begin{document}

{\small Akshat G \hfill \url{https://voyager2005.github.io} \par}
\vspace{4pt}
\hrule
\vspace{8pt}

My research interests lie in developing efficient, data-centric, and interpretable machine learning systems. I have always wanted to study this area in greater depth and pursue a career in research, as a research scientist developing machine learning systems that are efficient and interpretable. I am applying to the PhD program in Computer Science at The University of Texas at Austin as it provides an unparalleled opportunity to collaborate with experts, further my research, and lay the foundation for a fulfilling career in research.

\vspace{8pt}

\textbf{Research on Limiting Factors for a Model.} After my junior year, I found my research interests leaning toward understanding what actually limits a model's performance. I realized that rather than studying architectural variations in isolation, I wanted to understand when and why different designs worked better. In a benchmarking study of U-Net variants for lung segmentation~\hyperref[ref3]{[3]}, I compared U-Net, ResU-Net, Attention U-Net, and U-Net++ under identical conditions, finding that no single architecture was universally superior and that each was better suited to different segmentation demands. Under the guidance of Dr. Rashmi R, I tabulated these findings and presented them at the 10th Conference on Computer Vision and Image Processing (CVIP) at IIT Ropar.

This led me to question whether improving the data could matter more than improving the architecture. While working with the Kvasir-SEG benchmark, I found that its binary labels conflated surrounding tissue with camera artifacts, so I refined the ground truth into three classes and introduced preprocessing for specular reflection and contrast. A standard Attention U-Net trained on this corrected data reached a Dice score of 0.9525, a 1.2\% improvement over the strongest binary-task result I found in the literature. The result shifted my attention from architectural complexity toward a more fundamental question of how the data itself limits model performance.


\vspace{8pt}

\textbf{Research on Scarce and Unlabeled Data.} If data quality can matter more than architecture, the harder problem is when sufficient labeled data does not exist at all. I first approached this through student-teacher pseudo-labeling~\hyperref[ref5]{[5]}, where a teacher trained on labeled data guided a student using additional images generated by a GAN. However, the GAN's generations were not consistently reliable, allowing their errors to propagate into the student's training. This led me to develop A-PDF~\hyperref[ref6]{[6]}, a distribution-based framework that evaluated generated samples against real images and filtered unreliable ones before they reached the student. A-PDF reduced this problem, but it also exposed a deeper limitation of pseudo-labeling. The approach remained dependent on an intermediate model being sufficiently reliable, and filtering could not eliminate that dependency.

This pushed me toward diffusion pretraining instead. I explored generation as a pretext task for segmentation with the idea that in generating images, the model would inherently understand properties of the image like organ boundaries. I pretrained a Denoising Diffusion Probabilistic Model on unlabeled abdominal CT scans and transferred its learned structural representations to organ segmentation. This pretraining improved liver segmentation Dice from 0.75$\pm$0.36 to 0.93$\pm$0.16, while a frozen encoder, never fine-tuned on a single segmentation label, retained over 80\% of the fully fine-tuned model's performance. These results suggested that the pretraining phase was learning meaningful structural priors rather than superficial statistics. However, diffusion pretraining remains computationally expensive, and its transfer was noticeably stronger for some organs than others. These limitations have made me interested in understanding how structural priors can be learned more efficiently and why some learned priors transfer better across tasks than others.



\vspace{8pt}

\textbf{Research on Incorporating Reasoning into Architecture.} At Shell, I worked on detecting defects in gear systems, where a single gear can exhibit multiple defect types. Observing experienced inspectors, I noticed that they reasoned hierarchically, first identifying whether a defect was present, then narrowing its category before determining the specific defect and severity. Existing flat classifiers ignored this structure, producing predictions that could be difficult to reconcile across levels. I therefore built HiME, a Hierarchical Multi-label Estimator~\hyperref[ref8]{[8]}, which predicts each stage jointly while structurally enforcing consistency between them. This also made its predictions easier to interpret, as a human could trace the model's decision from broad defect category to specific diagnosis. On our deployment dataset, HiME reached 91\% accuracy with 100\% hierarchical consistency, against 88\% accuracy and 72\% consistency for a flat classifier.

\vspace{8pt}

These experiences have given me the questions I want to pursue, but not yet the depth to answer them rigorously. Graduate study at UT Austin would allow me to build that foundation. I see the PhD program in Computer Science at UT Austin as the next step toward developing the depth and independence I will need to pursue research as a career.

\vspace{8pt}

\textbf{Future Work. }The PhD program in Computer Science at UT Austin is particularly well suited to these interests because of its strength in generative modeling and efficient machine learning. I am especially interested in working with Prof. Qiang Liu, whose work on rectified flow~\hyperref[ref9]{[9]} and its extension to small, efficient one-step diffusion models~\hyperref[ref10]{[10]} examines exactly the cost problem I encountered in my own diffusion pretraining work. Prof. Mingyuan Zhou's work on score identity distillation~\hyperref[ref11]{[11]}, which distills pretrained diffusion models into fast, largely data-free one-step generators, connects closely with my interest in understanding how structural priors can be learned and transferred more efficiently, particularly the question of what can be preserved when a model is compressed this aggressively. I am also drawn to Prof. Haris Vikalo's group, whose work spans efficient inference under real resource constraints. After speaking with his student Haoran Zhang about his research on multi-layer hierarchical inference under partial and policy-dependent feedback, I was struck by how closely it parallels the problem I addressed in HiME, enforcing consistency across levels of a hierarchy, approached from the systems side rather than the classification side. I would be eager to explore this direction further within Prof. Vikalo's group and learn from the broader research community at UT Austin.

\vspace{4pt}
\hrule
\vspace{4pt}

{\footnotesize
\textbf{References}
\begin{enumerate}[leftmargin=1.3em, itemsep=1pt, label={[\arabic*]}]
\item \label{ref3} Akshat G, Divyansh Gupta, Rashmi R. Benchmarking U-Net Variants for Lung Segmentation in Chest X-Rays. \emph{Presented at the 10th Conference on Computer Vision and Image Processing (CVIP), IIT Ropar}, 2025. \url{https://voyager2005.github.io}
\item \label{ref4} Akshat G, Rashmi R. A Data-Centric Approach to Polyp Segmentation: Evaluating the Impact of Preprocessing and Ground-Truth Refinement. \emph{Under review, Scientific Reports}, 2026. \url{https://voyager2005.github.io}
\item \label{ref5} Akshat G, et al. Student-Teacher Pseudo-Labeling for Data-Scarce Medical Image Segmentation. \emph{Under Review, IEEE Access}, 2026. \url{https://voyager2005.github.io}
\item \label{ref6} Akshat G, et al. A-PDF: A Distribution-Based Evaluation Framework for Generative Data Augmentation Under Data Scarcity. \emph{Under Review IEEE Access}, 2026. \url{https://voyager2005.github.io}
\item \label{ref7} Akshat G, Divyansh Gupta, Shaleen Bhatnagar, Shilpa Ankalaki, Tusar Kanti Mishra. Unsupervised Anatomical Feature Learning via Diffusion Models: Enhanced Medical Image Segmentation with Denoising Diffusion Probabilistic Models. \emph{Under review, IEEE Access}, 2026. \url{https://voyager2005.github.io}
\item \label{ref8} Akshat Gururaj, Shashank Sharma, Ashish Mishra, Vishal R Ahuja. HiME: A Consistent, Encoder-Agnostic Hierarchical Multi-label Estimator. \emph{In Preparation for CVPR27}. \url{https://voyager2005.github.io}
\item \label{ref9} Xingchao Liu, Chengyue Gong, Qiang Liu. Flow Straight and Fast: Learning to Generate and Transfer Data with Rectified Flow. \emph{International Conference on Learning Representations (ICLR)}, 2023. \url{https://arxiv.org/abs/2209.03003}
\item \label{ref10} Yuanzhi Zhu, Xingchao Liu, Qiang Liu. SlimFlow: Training Smaller One-Step Diffusion Models with Rectified Flow. 2024. \url{https://arxiv.org/abs/2407.12718}
\item \label{ref11} Mingyuan Zhou, Huangjie Zheng, Zhendong Wang, Mingzhang Yin, Hai Huang. Score Identity Distillation: Exponentially Fast Distillation of Pretrained Diffusion Models for One-Step Generation. \emph{International Conference on Machine Learning (ICML)}, 2024. \url{https://arxiv.org/abs/2404.04057}
\end{enumerate}
}

\end{document}